% analysis/failure_modes.tex

\section{Failure Modes and Limitations of the AIM Pipeline}

% ============================================================
\subsection*{Motivation}
% ============================================================

Any modelling and optimization framework must account not only for its
capabilities, but also for its limitations.

Within the AIM pipeline, failure modes arise from the interaction of:
\begin{itemize}
\item coordinate systems,
\item geometric transformations,
\item physical realization,
\item numerical optimization.
\end{itemize}


% ============================================================
\subsection{CRS-Related Errors}
% ============================================================

Errors in coordinate reference systems include:
\begin{itemize}
\item incorrect CRS identification,
\item inconsistent transformations,
\item mixed coordinate systems within datasets.
\end{itemize}

\begin{proposition}
CRS inconsistencies lead to systematic spatial distortions that cannot
be corrected by local alignment optimization.
\end{proposition}


% ============================================================
\subsection{Projection Errors (World-to-Track)}
% ============================================================

The world-to-track transformation may fail due to:
\begin{itemize}
\item ambiguous nearest points,
\item high curvature regions,
\item sparse or noisy alignment representations.
\end{itemize}

\begin{proposition}
Projection errors introduce nonlocal distortions in the intrinsic
coordinate \(s\), affecting all downstream computations.
\end{proposition}


% ============================================================
\subsection{Model Mis-Specification}
% ============================================================

The parametric model may fail to represent the true alignment due to:
\begin{itemize}
\item insufficient transition families,
\item incorrect parameterization,
\item overly restrictive structure.
\end{itemize}

Hybrid correction terms mitigate but do not eliminate this issue.


% ============================================================
\subsection{Overfitting vs. Underfitting}
% ============================================================

Hybrid models introduce a trade-off:
\begin{itemize}
\item underfitting: parametric model too rigid,
\item overfitting: correction terms capture noise.
\end{itemize}

\begin{proposition}
Regularization is required to balance model fidelity and stability.
\end{proposition}


% ============================================================
\subsection{Realization Effects}
% ============================================================

The realization operator introduces:
\begin{itemize}
\item smoothing of curvature,
\item loss of high-frequency information,
\item dependence on construction processes.
\end{itemize}

\begin{proposition}
Certain features of the target alignment are not physically realizable
and will be systematically suppressed.
\end{proposition}


% ============================================================
\subsection{Inverse Problem Instability}
% ============================================================

The inverse realization problem is ill-posed:
\begin{itemize}
\item multiple target alignments yield similar realized geometry,
\item small changes in data may cause large parameter variations.
\end{itemize}

\begin{proposition}
Regularization and prior knowledge are essential for stable solutions.
\end{proposition}


% ============================================================
\subsection{Numerical Optimization Issues}
% ============================================================

Optimization may fail due to:
\begin{itemize}
\item poor initial guesses,
\item nonconvex objective functions,
\item ill-conditioned Jacobians.
\end{itemize}

Despite structured sparsity, convergence is not guaranteed.


% ============================================================
\subsection{Block Structure Degeneracy}
% ============================================================

The block structure may become degenerate if:
\begin{itemize}
\item coupling terms dominate,
\item insufficient data separates variables,
\item parameters become redundant.
\end{itemize}

\begin{proposition}
Loss of block separability reduces computational efficiency and
interpretability.
\end{proposition}


% ============================================================
\subsection{Measurement Noise and Bias}
% ============================================================

Measurement data may contain:
\begin{itemize}
\item random noise,
\item systematic bias,
\item missing or inconsistent segments.
\end{itemize}

\begin{proposition}
Noise propagates through the pipeline and may be amplified by inverse
operations.
\end{proposition}


% ============================================================
\subsection{Scale-Dependent Failures}
% ============================================================

Different failure modes dominate at different scales:
\begin{itemize}
\item large scale: CRS and projection errors,
\item medium scale: model mis-specification,
\item small scale: realization and measurement noise.
\end{itemize}


% ============================================================
\subsection{Non-Uniqueness of Solutions}
% ============================================================

Different alignment models may produce nearly identical realized
geometry.

\begin{proposition}
The mapping
\[
\kappa_{\mathrm{target}}
\;\mapsto\;
\gamma_{\mathrm{real}}
\]
is not injective.
\end{proposition}


% ============================================================
\subsection{Practical Engineering Limits}
% ============================================================

Certain effects are outside the scope of the model:
\begin{itemize}
\item material-dependent behaviour,
\item long-term degradation,
\item construction variability.
\end{itemize}

These factors must be treated as external constraints.


% ============================================================
\subsection{Key Insight}
% ============================================================

\begin{proposition}
The AIM pipeline is inherently approximate and limited by both physical
and computational constraints.
\end{proposition}


% ============================================================
\subsection{Connection to AIM}
% ============================================================

\begin{summary}
Failure modes in the AIM pipeline arise from the interaction of geometry,
physics, data, and computation.

Understanding these limitations is essential for:
\begin{itemize}
\item robust model design,
\item stable optimization,
\item reliable engineering application.
\end{itemize}

Thus, failure analysis is not an auxiliary component, but an integral
part of the AIM framework.
\end{summary}
